\begin{explanationbox}
	\textbf{\textcolor{CExplanation}{一句话直觉}}
	只要一条直线满足两个条件——它垂直于平面 $\alpha$，又在平面 $\beta$ 内——那么 $\alpha$ 和 $\beta$ 就互相垂直。

	\medskip
	\textbf{\textcolor{CExplanation}{核心拆解}}
	\begin{description}[style=nextline, leftmargin=2.8em, labelsep=0.5em, itemsep=0.45em]
		\item[\textbf{要点一（线面垂直）：}]
			直线 $l$ 必须与平面 $\alpha$ 垂直，即 $l\perp\alpha$。这意味着 $l$ 垂直于 $\alpha$ 内的任意直线。
		\item[\textbf{要点二（线在面内）：}]
			直线 $l$ 必须在平面 $\beta$ 内，即 $l\subset\beta$。这一条件不能省略。
	\end{description}

	\medskip
	\textbf{\textcolor{CExplanation}{几何本质（直二面角）}}
	当 $l\perp\alpha$ 且 $l\subset\beta$ 时，$l$ 垂直于 $\alpha$ 与 $\beta$ 的交线（因为交线在 $\alpha$ 内），从而对应的二面角的平面角为直角，故两平面垂直。

	\medskip
	\textbf{\textcolor{CExplanation}{代数意义（法向量垂直）}}
	设 $\alpha$ 的法向量为 $\boldsymbol{n}$，$\beta$ 的法向量为 $\boldsymbol{m}$，$l$ 的方向向量为 $\boldsymbol{a}$。由 $l\perp\alpha$ 得 $\boldsymbol{a}\parallel\boldsymbol{n}$；由 $l\subset\beta$ 得 $\boldsymbol{a}\perp\boldsymbol{m}$。于是 $\boldsymbol{n}\perp\boldsymbol{m}$，即 $\alpha\perp\beta$。

	\medskip
	\textbf{\textcolor{CExplanation}{考点价值（命题热点）}}
	立体几何解答题常在第一问考查面面垂直的证明。学生通常需要先寻找线面垂直关系，再转化为面面垂直。
	\begin{itemize}[leftmargin=*, itemsep=0.5em, topsep=0.4em]
		\item \textbf{考法一（直接应用）：} 题干直接给出线面垂直和线在面内条件，套用定理写出结论。
		\item \textbf{考法二（间接证明）：} 需要先通过线面垂直判定定理证明线面垂直，再得到面面垂直。
	\end{itemize}

	\medskip
	\textbf{\textcolor{CExplanation}{顿悟点}}
	判定面面垂直的关键是找到“桥梁线”——这条线既垂直于一个平面，又在另一个平面内。无需同时关心两个平面的垂直关系如何度量。

	\medskip
	\textbf{\textcolor{CExplanation}{使用场景}}
	\begin{itemize}[leftmargin=*, itemsep=0.5em, topsep=0.4em]
		\item \textbf{场景一（已知线面垂直）：} 在几何体中，若已知一条线段垂直于某个平面（如侧棱垂直于底面），且该线段位于某个侧面内，可直接判定面面垂直。
		\item \textbf{场景二（证明面面垂直）：} 当需要证明两个平面垂直时，首先考虑能否在其中一个平面内找到一条直线垂直于另一个平面。
	\end{itemize}

\end{explanationbox}
